\section*{Chapitre \ref{ch:g8:plastics} --- Les matières plastiques et les matériaux synthétiques}

\begin{solution}{exo:g8:plastics:1}
Naturels~: le coton, la laine. Artificiels~: le papier, la viscose.
Synthétiques~: le nylon, le polystyrène.
\end{solution}

\begin{solution}{exo:g8:plastics:2}
Du pétrole brut (ou du gaz naturel).
\end{solution}

\begin{solution}{exo:g8:plastics:3}
Le PET (polytéréphtalate d'éthylène)~: les bouteilles d'eau et de soda.
\end{solution}

\begin{solution}{exo:g8:plastics:4}
Un \omterm{def:g8:plastics:thermoplastic}{thermoplastique} se ramollit chaque fois qu'on le chauffe et peut être
moulé de nouveau~; un \omterm{def:g8:plastics:thermoplastic}{thermodurcissable} durcit pour de bon la première
fois et se carbonise au lieu de se ramollir.
\end{solution}

\begin{solution}{exo:g8:plastics:5}
Le polypropylène (PP), un \omterm{def:g8:plastics:thermoplastic}{thermoplastique}~: on peut le fondre et le
mouler de nouveau.
\end{solution}

\begin{solution}{exo:g8:plastics:6}
Le triangle ne fait que désigner la \omterm{def:g8:plastics:plastic}{matière plastique}, pour aider au
tri. Que l'objet soit recyclé dépend de sa collecte et de l'existence
d'une usine qui recycle cette matière.
\end{solution}

\begin{solution}{exo:g8:plastics:7}
Une poêle chauffe~: un \omterm{def:g8:plastics:thermoplastic}{thermoplastique} se ramollirait, un
\omterm{def:g8:plastics:thermoplastic}{thermodurcissable} ne se ramollit pas.
\end{solution}

\begin{solution}{exo:g8:plastics:8}
$\frac{353}{460} \approx 0.77$~: environ \qty{77}{\percent}.
\end{solution}

\begin{solution}{exo:g8:plastics:9}
Le PVC contient du chlore~: le \omterm{def:g4:what-a-fire-needs:burning}{brûler} dégage du chlorure d'hydrogène, un
gaz corrosif, en plus de la suie et du monoxyde de carbone dus à la
\omterm{def:g8:combustion-and-fuels:complete}{combustion incomplète} d'un feu de jardin.
\end{solution}

\begin{solution}{exo:g8:plastics:10}
$300 \times 36 = 10\,800$ bouteilles~; $10\,800 \times 25 = 270\,000$\,g,
soit \qty{270}{kg}.
\end{solution}

\begin{solution}{exo:g8:plastics:11}
Du dioxyde de carbone et de l'eau. Son carbone vient du pétrole brut,
stocké sous terre depuis des millions d'années~: le \omterm{def:g4:what-a-fire-needs:burning}{brûler} ajoute du
dioxyde de carbone à l'\omterm{def:g4:air-a-mixture-of-gases:air}{air}, comme le fait la \omterm{def:g4:what-a-fire-needs:burning}{combustion} d'un \omterm{def:g8:combustion-and-fuels:fossil}{combustible
fossile}.
\end{solution}

\begin{solution}{exo:g8:plastics:12}
La bouteille ne pourrit pas~: le soleil et les vagues la brisent en
morceaux de plus en plus petits. Les minuscules morceaux flottent ou
coulent dans l'eau, et les poissons les avalent avec leur nourriture.
\end{solution}

\begin{solution}{pb:g8:plastics:1}
\textbf{1.} Corps~: 1 (PET). Bouchons~: 2 (PEHD) ou 5 (PP).

\textbf{2.} Synthétiques~: elles sont fabriquées à partir de substances
construites par les chimistes, surtout à partir du pétrole.

\textbf{3.} Oui~: le PET est un \omterm{def:g8:plastics:thermoplastic}{thermoplastique}.

\textbf{4.} Chaque \omterm{def:g8:plastics:plastic}{matière plastique} se recycle séparément~: fondu
ensemble, un \omterm{def:g6:pure-substances-and-mixtures:mixture}{mélange} de \omterm{def:g8:plastics:plastic}{matières plastiques} donne un matériau de
mauvaise qualité.

\textbf{5.} $0.80 \times 1000 = \qty{800}{kg}$.

\textbf{6.} Bouchons~: $0.12 \times 1000 = \qty{120}{kg}$. Étiquettes~: $0.05
\times 1000 = \qty{50}{kg}$.

\textbf{7.} $0.03 \times 1000 = \qty{30}{kg}$ de saletés et de liquide.

\textbf{8.} $80 + 12 + 5 = \qty{97}{\percent}$.

\textbf{9.} $0.025 \times 800 = \qty{20}{kg}$.

\textbf{10.} Les paillettes économisent du pétrole brut et l'énergie
nécessaire pour fabriquer du PET neuf, et elles transforment un déchet
en objet utile.

\textbf{11.} Du dioxyde de carbone (et, si la \omterm{def:g4:what-a-fire-needs:burning}{combustion} est mauvaise, du
monoxyde de carbone et de la suie).

\textbf{12.} $800 - 20 = \qty{780}{kg}$ de paillettes de PET propres par
balle.
\end{solution}
